\documentclass{subfiles}
\begin{document}
    Our discussion on self-adjusting lists will proceed as follow: 
        first, we introduce the definition of self-adjusting list;
        then, we focus our attention on some of the most common rules used to keep a list organized.

    Conceptually, self-adjusting lists are trivial:
        these are lists in which, after accessing or inserting an item, 
        some update procedure is called to keep the list organized.
        These procedure are defined according to some rules,
        among which those that stands out are 
        \begin{itemize}
            \item \emph{\gls{fc}}: we mantain an array that stores at each position \(i\),
                the count of how many times item \(i\) has been accessed or added;
            \item \emph{\gls{mtf}:} each time an item is added, it is moved to the head of the list,
                while each accessed item is moved slighly closer to the root;
            \item \emph{\gls{t}:} accessing or inserting an item, implies an exchange of the added/%
                accessed item with the one on its left.
        \end{itemize}
    
    Let \(\E{A}\) be the expected cost to access an item using algorithm \(A\), 
        where \(p = \set{p_{1}, p_{2}, \ldots, p_{n}}\) is the probability distribution associated
        to the sequence of operations. Assume \(DP\) to be an algorithm that 
        orders the elements according to \(p\). By the law of large numbers it follows \(\E{FC}/\E{DP} = 1\).
        Additionally, one could prove that \(\E{MTF} \le 2 \E{DP}\),
        and as proved in \cite{Rivest1976} \(\E{T} \le \E{MTF}\). 
        A similar result can be easily derived for insertions and deletions.
        From the above, it follows that, at least in theory, 
        both FC and transpose out perform MTF. As we are about to show, theory and practice not always coincide.
        What we want to show is that, although in theory FC and transpose are better then MTF,
        in practical applications MTF performs better.
        
    For any algorithm \(A\), let us denote with \(C_{A}(s)\) the total cost of all operations,
        by \(F_{A}(s)\) the number of exchanges that move an item to the head of the list,
        and by \(P_{A}(s)\) any other exchange. The following holds.

    \begin{theorem}[\cite[Theorem 1.]{SleatorT1985}]\label[theorem]{thm:SleatorT1985-thm-1}
        For any Algorithm \(A\) and any sequence of operations \(s\) starting with the empty set
        \[
            C_{MTF}(s) \le C_{A}(s) + P_{A}(s) - F_{A}(s) - m.
        \]
    \end{theorem}
    \Cref{thm:SleatorT1985-thm-1} and its generalized version (see \cite[Theorem 2.]{SleatorT1985}),
    show that MTF differs by any other algorithm by a most a factor of 2. No analogous result 
    holds for FC and transpose.
\end{document}
